\begin{table}[!htbp]
\centering
\begin{threeparttable}
\caption{Distribution of primary EB-shrunken school value added}
\label{tab:school-va-distribution}
\footnotesize
\setlength{\tabcolsep}{2pt}
\begin{tabular}{lcc|cccccc}
\toprule
 & \multicolumn{2}{c|}{Outcome distribution} & \multicolumn{6}{c}{Value-added distribution} \\
\cmidrule(lr){2-3} \cmidrule(lr){4-9}
Outcome & \shortstack{Sample\\Mean} & \shortstack{Sample\\SD} & VA SD & P10 & P25 & P50 & P75 & P90 \\
\midrule
Math & 0.043 & 1.002 & 0.222 & -0.195 & -0.141 & -0.062 & 0.057 & 0.297 \\
Verbal & 0.049 & 0.989 & 0.150 & -0.172 & -0.111 & -0.021 & 0.083 & 0.202 \\
High-premium institution & 0.142 & 0.349 & 0.106 & -0.090 & -0.066 & -0.029 & 0.023 & 0.144 \\
High-premium field & 0.347 & 0.476 & 0.053 & -0.060 & -0.035 & -0.004 & 0.030 & 0.063 \\
Log proj. income & 13.893 & 0.522 & 0.072 & -0.093 & -0.049 & -0.002 & 0.044 & 0.090 \\
\bottomrule
\end{tabular}
\begin{tablenotes}[flushleft]
\footnotesize
\item[] \parbox{0.84\textwidth}{Notes: The table reports student-weighted distributions of EB-shrunken, student-centered school value added. Schools are weighted by the number of students in the corresponding Stata value-added regression. Sample moments refer to the outcome-specific value-added regression sample. Math and verbal are measured in admission-test standard deviations; the two high-return outcomes are binary; projected program income is measured in logs.}
\end{tablenotes}
\end{threeparttable}
\end{table}
